% Options for packages loaded elsewhere
\PassOptionsToPackage{unicode}{hyperref}
\PassOptionsToPackage{hyphens}{url}
\documentclass[
]{article}
\usepackage{xcolor}
\usepackage[margin=1in]{geometry}
\usepackage{amsmath,amssymb}
\setcounter{secnumdepth}{-\maxdimen} % remove section numbering
\usepackage{iftex}
\ifPDFTeX
  \usepackage[T1]{fontenc}
  \usepackage[utf8]{inputenc}
  \usepackage{textcomp} % provide euro and other symbols
\else % if luatex or xetex
  \usepackage{unicode-math} % this also loads fontspec
  \defaultfontfeatures{Scale=MatchLowercase}
  \defaultfontfeatures[\rmfamily]{Ligatures=TeX,Scale=1}
\fi
\usepackage{lmodern}
\ifPDFTeX\else
  % xetex/luatex font selection
\fi
% Use upquote if available, for straight quotes in verbatim environments
\IfFileExists{upquote.sty}{\usepackage{upquote}}{}
\IfFileExists{microtype.sty}{% use microtype if available
  \usepackage[]{microtype}
  \UseMicrotypeSet[protrusion]{basicmath} % disable protrusion for tt fonts
}{}
\makeatletter
\@ifundefined{KOMAClassName}{% if non-KOMA class
  \IfFileExists{parskip.sty}{%
    \usepackage{parskip}
  }{% else
    \setlength{\parindent}{0pt}
    \setlength{\parskip}{6pt plus 2pt minus 1pt}}
}{% if KOMA class
  \KOMAoptions{parskip=half}}
\makeatother
\usepackage{color}
\usepackage{fancyvrb}
\newcommand{\VerbBar}{|}
\newcommand{\VERB}{\Verb[commandchars=\\\{\}]}
\DefineVerbatimEnvironment{Highlighting}{Verbatim}{commandchars=\\\{\}}
% Add ',fontsize=\small' for more characters per line
\usepackage{framed}
\definecolor{shadecolor}{RGB}{248,248,248}
\newenvironment{Shaded}{\begin{snugshade}}{\end{snugshade}}
\newcommand{\AlertTok}[1]{\textcolor[rgb]{0.94,0.16,0.16}{#1}}
\newcommand{\AnnotationTok}[1]{\textcolor[rgb]{0.56,0.35,0.01}{\textbf{\textit{#1}}}}
\newcommand{\AttributeTok}[1]{\textcolor[rgb]{0.13,0.29,0.53}{#1}}
\newcommand{\BaseNTok}[1]{\textcolor[rgb]{0.00,0.00,0.81}{#1}}
\newcommand{\BuiltInTok}[1]{#1}
\newcommand{\CharTok}[1]{\textcolor[rgb]{0.31,0.60,0.02}{#1}}
\newcommand{\CommentTok}[1]{\textcolor[rgb]{0.56,0.35,0.01}{\textit{#1}}}
\newcommand{\CommentVarTok}[1]{\textcolor[rgb]{0.56,0.35,0.01}{\textbf{\textit{#1}}}}
\newcommand{\ConstantTok}[1]{\textcolor[rgb]{0.56,0.35,0.01}{#1}}
\newcommand{\ControlFlowTok}[1]{\textcolor[rgb]{0.13,0.29,0.53}{\textbf{#1}}}
\newcommand{\DataTypeTok}[1]{\textcolor[rgb]{0.13,0.29,0.53}{#1}}
\newcommand{\DecValTok}[1]{\textcolor[rgb]{0.00,0.00,0.81}{#1}}
\newcommand{\DocumentationTok}[1]{\textcolor[rgb]{0.56,0.35,0.01}{\textbf{\textit{#1}}}}
\newcommand{\ErrorTok}[1]{\textcolor[rgb]{0.64,0.00,0.00}{\textbf{#1}}}
\newcommand{\ExtensionTok}[1]{#1}
\newcommand{\FloatTok}[1]{\textcolor[rgb]{0.00,0.00,0.81}{#1}}
\newcommand{\FunctionTok}[1]{\textcolor[rgb]{0.13,0.29,0.53}{\textbf{#1}}}
\newcommand{\ImportTok}[1]{#1}
\newcommand{\InformationTok}[1]{\textcolor[rgb]{0.56,0.35,0.01}{\textbf{\textit{#1}}}}
\newcommand{\KeywordTok}[1]{\textcolor[rgb]{0.13,0.29,0.53}{\textbf{#1}}}
\newcommand{\NormalTok}[1]{#1}
\newcommand{\OperatorTok}[1]{\textcolor[rgb]{0.81,0.36,0.00}{\textbf{#1}}}
\newcommand{\OtherTok}[1]{\textcolor[rgb]{0.56,0.35,0.01}{#1}}
\newcommand{\PreprocessorTok}[1]{\textcolor[rgb]{0.56,0.35,0.01}{\textit{#1}}}
\newcommand{\RegionMarkerTok}[1]{#1}
\newcommand{\SpecialCharTok}[1]{\textcolor[rgb]{0.81,0.36,0.00}{\textbf{#1}}}
\newcommand{\SpecialStringTok}[1]{\textcolor[rgb]{0.31,0.60,0.02}{#1}}
\newcommand{\StringTok}[1]{\textcolor[rgb]{0.31,0.60,0.02}{#1}}
\newcommand{\VariableTok}[1]{\textcolor[rgb]{0.00,0.00,0.00}{#1}}
\newcommand{\VerbatimStringTok}[1]{\textcolor[rgb]{0.31,0.60,0.02}{#1}}
\newcommand{\WarningTok}[1]{\textcolor[rgb]{0.56,0.35,0.01}{\textbf{\textit{#1}}}}
\usepackage{graphicx}
\makeatletter
\newsavebox\pandoc@box
\newcommand*\pandocbounded[1]{% scales image to fit in text height/width
  \sbox\pandoc@box{#1}%
  \Gscale@div\@tempa{\textheight}{\dimexpr\ht\pandoc@box+\dp\pandoc@box\relax}%
  \Gscale@div\@tempb{\linewidth}{\wd\pandoc@box}%
  \ifdim\@tempb\p@<\@tempa\p@\let\@tempa\@tempb\fi% select the smaller of both
  \ifdim\@tempa\p@<\p@\scalebox{\@tempa}{\usebox\pandoc@box}%
  \else\usebox{\pandoc@box}%
  \fi%
}
% Set default figure placement to htbp
\def\fps@figure{htbp}
\makeatother
\setlength{\emergencystretch}{3em} % prevent overfull lines
\providecommand{\tightlist}{%
  \setlength{\itemsep}{0pt}\setlength{\parskip}{0pt}}
\usepackage{bookmark}
\IfFileExists{xurl.sty}{\usepackage{xurl}}{} % add URL line breaks if available
\urlstyle{same}
\hypersetup{
  pdftitle={PEC5: Contraste de hipótesis},
  pdfauthor={Ivan Miranda Moral},
  hidelinks,
  pdfcreator={LaTeX via pandoc}}

\title{PEC5: Contraste de hipótesis}
\author{Ivan Miranda Moral}
\date{2026-05-08}

\begin{document}
\maketitle

\subsection{Introducción}\label{introducciuxf3n}

El objetivo de esta PEC es analizar la gestión de residuos en una
muestra de municipios. Los datos se encuentran en el archivo
\emph{Municipios\_Sostenibles.csv}. Este dataset incluye información
sobre la población, el tipo de municipio y los principales indicadores
de reciclaje.

\subsubsection{Variables:}\label{variables}

\begin{itemize}
\tightlist
\item
  \emph{Municipio\_ID}: Identificador único.
\item
  \emph{Habitantes}: Número de personas residentes.
\item
  \emph{Tipo\_Municipio}: Categoría (Urbano, Rural, Turístico).
\item
  \emph{Residuos\_kg\_capita}: Kilogramos de residuos generados por
  persona y año.
\item
  \emph{Tasa\_Reciclaje}: Porcentaje de recogida selectiva.
\item
  \emph{Coste\_Gestion\_EUR}: Coste por habitante de la gestión de
  residuos.
\item
  \emph{Puerta\_a\_Puerta}: Si el municipio tiene implementado este
  sistema (Sí/No).
\end{itemize}

\begin{Shaded}
\begin{Highlighting}[]
\NormalTok{datos }\OtherTok{\textless{}{-}} \FunctionTok{read.csv}\NormalTok{(}\StringTok{"Municipios\_Sostenibles 2.csv"}\NormalTok{,}
                  \AttributeTok{header =} \ConstantTok{TRUE}\NormalTok{, }\AttributeTok{encoding =} \StringTok{"UTF{-}8"}\NormalTok{,}
                  \AttributeTok{stringsAsFactors =} \ConstantTok{TRUE}\NormalTok{)}
\end{Highlighting}
\end{Shaded}

\subsection{Problema 1}\label{problema-1}

(3,5 puntos)La agencia ambiental sospecha que, debido a nuevas
políticas, la media poblacional de la \emph{Tasa\_Reciclaje} en el
conjunto de los municipios es superior al 38\%.

\begin{enumerate}
\def\labelenumi{\alph{enumi})}
\tightlist
\item
  (1,5 puntos) Formule el contraste de hipótesis adecuado para verificar
  esta sospecha. Realice el contraste utilizando un nivel de
  significación del 5\% (𝛼 = 0.05), empleando exclusivamente la
  instrucción directa de R. Razone la conclusión basándose en el p-valor
  y compruebe si el resultado es coherente con el intervalo de confianza
  proporcionado por la salida del programa.
\end{enumerate}

\begin{quote}
\textcolor{blue}{El conjunto \texttt{datos\_activos}
filtra los registros en los que el uso de CPU es superior al 10\%, con el objetivo de analizar únicamente situaciones de carga real del sistema.}
\end{quote}

\begin{quote}
\textcolor{blue}{ 
A continuación, se muestran las primeras 6 observaciones del nuevo conjunto de datos:
}
\end{quote}

\begin{enumerate}
\def\labelenumi{\alph{enumi})}
\setcounter{enumi}{1}
\tightlist
\item
  (2 puntos) Realice el mismo contraste utilizando las fórmulas de los
  módulos. Debe indicar:
\end{enumerate}

\begin{enumerate}
\def\labelenumi{\arabic{enumi}.}
\tightlist
\item
  Hipótesis nula (𝐻0) y alternativa (𝐻1) utilizando la simbología
  correcta (𝜇, 𝜇0).
\item
  Cálculo del estadístico de contraste mediante la fórmula
  correspondiente.
\item
  Cálculo del p-valor mediante R.
\item
  Decisión y conclusión final aplicada al contexto del problema.
\end{enumerate}

\begin{quote}
\textcolor{blue}{ 
Finalmente
}
\end{quote}

\begin{quote}
Referencias:
\begin{itemize}
\item \textbf{Análisis de datos y estadística descriptiva con R}, unidad \textbf{3. Análisis demográfico en Cataluña}, apartado \textbf{3.1. Manejo de conjuntos de datos}, p. 31.
\end{itemize}
\end{quote}

\subsection{Problema 2}\label{problema-2}

(3 puntos)Se quiere determinar si la proporción de municipios que han
implementado el sistema \emph{Puerta\_a\_Puerta} es igual al 50\% o si
es diferente. Utilice un nivel de significación del 5\% (𝛼 = 0.05) para
realizar el contraste de hipótesis. Debe indicar:

\begin{enumerate}
\def\labelenumi{\arabic{enumi}.}
\tightlist
\item
  El tipo de contraste (unilateral o bilateral) y las hipótesis 𝐻0 y 𝐻1.
\item
  Realizar el test correspondiente con R.
\item
  La interpretación del p-valor obtenido para llegar a la conclusión.
\end{enumerate}

\begin{quote}
\textcolor{blue}{
Se ha calculado un intervalo de confianza del 95\% para la media poblacional de la latencia utilizando la instrucción:
}
\end{quote}

\begin{quote}
Referencias:
\begin{itemize}
\item \textbf{Intervalos de confianza}, unidad \textbf{Introducción a los intervalos de confianza. El caso de la media aritmética}, apartado \textbf{3. Intervalo de confianza para la media cuando la población es normal y desconocemos la desviación típica}, p. 11--13.
\item \textbf{Teorema del límite central}, unidad \textbf{La distribución de la media muestral}, apartado \textbf{1.2. Caso de desviación típica poblacional desconocida. La t de Student}, p. 8--10.
\end{itemize}
\end{quote}

\subsection{Problema 3}\label{problema-3}

(3,5 puntos) Queremos evaluar si existe una diferencia significativa
entre el \emph{Coste\_Gestion\_EUR} de los municipios de tipo ``Urbano''
frente a los de tipo ``Rural''. Asumiendo que las varianzas
poblacionales son desconocidas pero iguales, realice el contraste con un
nivel de significación del 1\% (𝛼 = 0.01).)

\begin{enumerate}
\def\labelenumi{\arabic{enumi}.}
\tightlist
\item
  Defina las hipótesis nula y alternativa.
\item
  Prepare los subconjuntos de datos en R para ambos tipos de municipio.
\item
  Ejecute el test de contraste correspondiente,
\item
  Indique el p-valor, la decisión tomada y calcule el intervalo de
  confianza para la diferencia de medias, justificando la relación entre
  este intervalo y la conclusión final.
\end{enumerate}

\begin{quote}
\textcolor{blue}{
Queremos que el intervalo de confianza del 95\% para la media tenga un error máximo de 0.2 ms.
}
\end{quote}

\begin{quote}
\textcolor{blue}{
En el problema anterior se trabajó con un intervalo de confianza para la media poblacional con varianza desconocida, por lo que el valor crítico sigue una distribución \(t\) de Student y el intervalo tiene la forma:
}
\end{quote}

\begin{quote}
Referencias:
\begin{itemize}
\item \textbf{Intervalos de confianza}, unidad \textbf{Introducción a los intervalos de confianza. El caso de la media aritmética}, apartado \textbf{3. Intervalo de confianza para la media cuando la población es normal y desconocemos la desviación típica}, p. 11--13.
\item \textbf{Intervalos de confianza}, unidad \textbf{Introducción a los intervalos de confianza. El caso de la media aritmética}, apartado \textbf{3.1. El tamaño de la muestra}, p. 13--14.
\item \textbf{Distribuciones de probabilidad e inferencia estadística con R}, unidad \textbf{2. Inferencia estadística}, apartado \textbf{2.2. Inferencia sobre la media con varianza poblacional desconocida}, p. 27.
\end{itemize}
\end{quote}

\end{document}
